\section{绪论}

去中心化金融的兴起是自电子交易出现以来金融中介领域最重大的结构性变革之一。自2020年1月至2025年12月间，DeFi协议的总锁仓价值（TVL）从未及10亿美元扩展至峰值超过1000亿美元，其中借贷协议始终占据最大份额~\cite{schar2021decentralized2, jiang2023defi}。传统金融中介依赖机构信用评估、关系型贷款和监管合规框架来管理风险，而去中心化金融借贷协议则通过不可篡改的智能合约运行，由算法自动执行抵押品要求、利率调整以及抵押不足头寸的自动清算~\cite{gudgeon2020defi}。这些协议的完全透明性意味着每一笔存款、取款、借款、还款、抵押品调整和清算事件都不可更改地记录在公共区块链上，为研究金融市场参与者在不同程度风险敞口下的行为提供了前所未有的经验观察窗口。

这种透明性催生了一笔从多个分析角度考察DeFi借贷的快速扩展的学术文献。第一类文献研究这些协议固有的系统性风险，包括流动性危机、连锁清算以及可组合DeFi原语之间的传染效应~\cite{qin2021empirical, perez2020liquidations}。第二类文献开发了适合DeFi借贷独特特征的风险建模框架，涵盖多链环境、跨协议依赖关系以及机构级风险分类体系~\cite{oberholzer2026toward, sevim2026interoperability}。第三类文献将机器学习方法应用于链上数据进行信用评分和违约预测~\cite{ghosh2024onchain}。第四类文献则研究具体协议机制的工程特性，包括清算阈值、利率曲线和治理设计~\cite{iftikhar2025automated, tovanich2023contagion}。

尽管已有如此丰富的文献，其核心处仍存在一个显著的研究空白。现有关于DeFi借贷风险的研究绝大多数聚焦于可称为"状态变量"的指标——由健康因子、抵押率以及离散时间点上借款金额所描绘的头寸健康状况静态快照。对现有文献而言，鲜有得到系统经验研究的是借款人应对头寸风险不断上升所经历的\textbf{行为过程}：即在风险状况恶化时，采取主动决策调整抵押品、偿还债务、增加敞口或维持头寸。这种区分的意义在于，两个在特定时刻具有相同健康因子的借款人，可能因为应对不断累积的风险压力时的行为响应方式不同而遵循了本质上不同的轨迹。一个借款人可能在风险周期早期主动追加抵押品并降低杠杆，而另一个则可能维持甚至增加风险敞口，直至被第三方套利者强制平仓。传统的风险指标会将这两个借款人置于相同的风险水平，但其行为过程暗示两者在风险管理方式上存在对后续结果可能具有信息含量的显著差异。

本研究报告通过在DeFi借贷基础设施、行为金融理论和链上信用风险评估三者的交叉点上提出了填补上述空白的研究设计。我们提出的核心研究问题是：\textbf{在公开链上借贷市场中，借款人在头寸风险上升期间所采取的主动调整行为，是否提供了超越常规链上指标所包含的风险信息的增量信息？}这一问题可分解为两个分析上相互独立但在经验上紧密相关的层面。第一个行为层面询问：当贷款头寸接近风险阈值时，借款人是否表现出系统的、非随机的主动调整模式，这些模式可以区别于价格变动的机械效应和清算机器人的操作？第二个信用导向层面则询问：即使控制了包括当前健康因子、抵押率、账户活动、资产集中度和市场波动性在内的常规链上风险指标后，能否从捕捉主动调整路径的\textbf{行为过程变量}中获得对未来清算事件或风险恶化和统计经济上显著的增量预测能力？

本报告结构如下：第二节在三个领域中普查了相关文献：DeFi借贷基础设施与风险、行为金融与前景理论以及链上信用风险评估。第三节明确了具体的研究课题，并阐述了研究问题和可检验假设。第四节描述了所提出的研究方法论，包括数据来源、变量构建和分析框架。第五节讨论了可能的研究成果、贡献、局限性及边界条件。

%
% Figure 1: Conceptual Framework
%
\begin{figure}[t]
\centering
\resizebox{\textwidth}{!}{%
\begin{tikzpicture}[
  node distance=1.2cm and 0.6cm,
  box/.style={rectangle, rounded corners=2pt, draw=#1!60, fill=#1!12, text=#1!80!black, minimum width=2.2cm, minimum height=0.65cm, align=center, font=\scriptsize\sffamily, line width=0.4pt},
  widebox/.style={rectangle, rounded corners=2pt, draw=#1!60, fill=#1!12, text=#1!80!black, minimum width=9cm, minimum height=0.65cm, align=center, font=\small\sffamily\bfseries, line width=0.4pt},
  smallbox/.style={rectangle, rounded corners=2pt, draw=#1!60, fill=#1!10, text=#1!80!black, minimum width=1.8cm, minimum height=0.5cm, align=center, font=\tiny\sffamily, line width=0.3pt},
  arrow/.style={->, >=stealth, line width=0.6pt, draw=black!55},
  darrow/.style={<->, >=stealth, line width=0.7pt, draw=black!55},
]
  % Layer 1: Protocols
  \node[widebox=blue, label={[font=\footnotesize\sffamily\bfseries]north:Data Sources}] (proto) {Aave (V2/V3) \hspace{0.8cm} Compound (V2/V3) \hspace{0.8cm} MakerDAO};

  % Layer 2: Behavioral variables  
  \node[widebox=teal, below=0.7cm of proto, label={[font=\footnotesize\sffamily\bfseries]north:Behavioral Process Extraction}] (behave) {\normalsize Borrower Active Adjustment \quad (collateral changes, debt management, response timing)};

  % Layer 3: Two tracks
  \node[box=violet, below left=1.1cm and -1.2cm of behave, minimum width=3.8cm, font=\scriptsize\sffamily] (rq1) {\textbf{RQ1: Behavioral Layer}\\[2pt]Do borrowers systematically adjust\\as HF $\rightarrow$ 1.0?};
  \node[box=violet, below right=1.1cm and -1.2cm of behave, minimum width=3.8cm, font=\scriptsize\sffamily] (rq2) {\textbf{RQ2: Credit Layer}\\[2pt]Does adjustment behavior predict\\future liquidations?};

  \node[box=red!80!black, below=0.6cm of rq1, minimum width=3.8cm, font=\scriptsize\sffamily] (out1) {Prospect Theory Validation\\(loss aversion, reference-point)};
  \node[box=red!80!black, below=0.6cm of rq2, minimum width=3.8cm, font=\scriptsize\sffamily] (out2) {Behavioral Early Warning\\(credit risk signals)};

  % Arrows
  \draw[arrow, thick] (proto.south) -- (behave.north);
  \draw[arrow, thick] (behave.south) |- (rq1.north);
  \draw[arrow, thick] (behave.south) |- (rq2.north);
  \draw[arrow] (rq1.south) -- (out1.north);
  \draw[arrow] (rq2.south) -- (out2.north);
  \draw[darrow, thick] (rq1.east) -- node[above, font=\tiny\sffamily]{H2a: Incremental predictive power?} (rq2.west);

  % Left labels
  \node[font=\footnotesize\sffamily, text=black!60, rotate=90, above left=0.0cm and -1.8cm of behave, anchor=south] {$\longleftarrow$ Behavioral Finance $\longrightarrow$};
\end{tikzpicture}%
}
\caption{Conceptual framework for studying borrower behavior and credit signals in DeFi lending. The framework proceeds from raw on-chain data (top) through behavioral process variable extraction (center) to the research questions (bottom). RQ1 characterizes how borrowers adjust positions as risk accumulates; RQ2 tests whether these behavioral patterns provide incremental predictive power for future liquidation events beyond conventional risk indicators.}
\label{fig:framework}
\end{figure}

\vspace{4pt}
